\documentclass[aspectratio=169]{beamer}
\usepackage[norsk]{babel}
\usepackage{booktabs}
\usepackage[T1]{fontenc}
\usepackage[utf8]{inputenc}
\usepackage{graphicx,epstopdf,tikz}
\usepackage{siunitx, enumerate, xcolor}
\usepackage{ enumerate, xcolor}
\usepackage[export]{adjustbox}
\usepackage{circuitikz}
%\usepackage{enumitem}
%\usepackage{enumitem, siunitx}
\usepackage{multicol}
%\usepackage{mathtools}
\usetikzlibrary{decorations.pathreplacing, calligraphy, calc}

\usepackage{minted}

\usetheme[slogan=norsk,style=vertical, mathfont=serif]{NTNU}
\setbeamertemplate{section in toc}[circle]

\title[rl]{Programmering og data i økonomiske fag}
\date{22. August 2024}
\author{Jonas J. Harang}

%SIunitx
\let\DeclareUSUnit\DeclareSIUnit
\let\US\SI
\let\us\si

\DeclareUSUnit\gallon{gal}
\DeclareUSUnit\inch{in}
\DeclareUSUnit\psi{psi}

\DeclareSIUnit\atm{atm}
\DeclareSIUnit\bar{bar}
\DeclareSIUnit\mmhg{mm Hg}

%%--- tikzcolor for ntnu
\definecolor{NTNUblue}{HTML}{00509e}
\definecolor{NTNUlightblue}{HTML}{6096d0}
\definecolor{NTNUorange}{HTML}{ef8114}
\definecolor{NTNUbrown}{HTML}{cfb887}
\definecolor{NTNUred}{HTML}{b01b81}
\definecolor{NTNUgrey}{HTML}{bebebe}
\definecolor{NTNUcyan}{HTML}{3cbfbe}
\definecolor{NTNUviolet}{HTML}{482776}



%-----------Beamer jonas
%\setbeamertemplate{subsection in toc}[subsections numbered]
\setbeamertemplate{subsection in toc}[subsections unnumbered]
\makeatletter
\newcommand{\SetSectionNumber}[1]{%
\beamer@tocsectionnumber=\numexpr#1-1}
\makeatother



\begin{document}

\SetSectionNumber{2}
%
\begin{frame}
   \maketitle
\end{frame}

\begin{frame}{}
  \begin{columns}
    \begin{column}{0.5\textwidth}
      \begin{block}{IIRA2001:}
        \begin{itemize}
          \item {Grunnleggende programmering i Python}
          \item {Python som verktøy i økonomiske fag}
        \end{itemize}
      \end{block}
    \end{column}
    \begin{column}{0.5\textwidth}
      \onslide<2->{Faglig innhold utvikles i samarbeid mellom IIR og IIF og fokuset er på å 
      sette programmering inn i fagene dere tar, og gi dere et kraftig verktøyssett til 
      problemene dere jobber med under studiene}
    \end{column}
  \end{columns}
\end{frame}

\begin{frame}{Hvordan?}
  \begin{itemize}
    \item {2x4 timers læringsøkter i storhavet (255A/B) på NMK tors/fre}
    \item {Øktene er fleksible med tanke på innhold men utgangspunktet er:}
        \item {Første 2 timer: Studentaktiv forelesning hvor vi går igjennom programmeringskonsepter sammen}
        \item {Siste 2 timer: Arbeid med oppgaver, under veiledning av foreleser og/eller studentassistenter}
  \end{itemize}
\end{frame}

\begin{frame}{Pensum}
  \begin{columns}
    \begin{column}{0.5\textwidth}
      Ingen pensumbok:
        \begin{itemize}
          \item {Relevante ressurser gjøres tilgjengelig på blackboard}
          \item {Eller lenkes til i forelesningsmateriell}
          \item {Tidligere studenter har hatt litt nytte av \href{https://www.akademika.no/teknologi/data-og-informasjonsteknologi/python-realfag/9788245046052?gad_source=1&gclid=EAIaIQobChMIh5HM4fqFiAMVH1iRBR18gQPXEAAYAiAAEgI-E_D_BwE}{Python for realfag}}
        \end{itemize}
    \end{column}

    \begin{column}{0.5\textwidth}
      Inneholdet i forelesning og øvinger/oppgaver er pensum
    \end{column}
  \end{columns}
\end{frame}

\begin{frame}{Øvinger}
  \begin{columns}
    \begin{column}{0.5\textwidth}
      \begin{block}{Større øvingsoppgaver:}
        \begin{itemize}
          \item {I hovedsak ukentlige øvingsoppgaver}
          \item {Noen øvingsoppgaver vil jobbes med over flere uker}
          \item {Øvingsoppgavene vil være nært knyttet til \textbf{mappeinnleveringen}}
          \item {Ikke obligatoriske (i praksis må de jo gjøres)}
        \end{itemize}
      \end{block}
    \end{column}
    \begin{column}{0.5\textwidth}
      \onslide<2->{\begin{itemize}
        \item {Noen selvrettende øvinger i Moodle/CodeRunner}
        \item {Korte python-oppgaver for å øve på det tekniske}
        \item {\color{gray}{De automatisk rettede øvingene er obligatoriske}}
        \item {De er å finne \href{iirmoodle.it.ntnu.no}{her}}
           \begin{itemize}
              \item {Logg på med Feide og meld deg opp i kurset}
           \end{itemize}
      \end{itemize}}
    \end{column}
  \end{columns}
\end{frame}

\begin{frame}{Vurderingsform}
  \begin{columns}
    \begin{column}{0.5\textwidth}
      \begin{block}{Mappevurdering}
        \begin{itemize}
          \item {Ingen skriftlig avsluttende skoleeksamen: dere skal levere en \textbf{mappe} med arbeider til vurdering}
          \item {Mappen skal i utgangspunktet ha 2 deler}
          \item {Del 1 vil være nært knyttet til øvingsoppgavene}
          \item {Del 2 vil være meget åpen og knyttet til data-analyse}
          \item {Oppgavene skal gi rom for kreativitet og selvstendighet, samt fleksibilitet for de som sliter}
        \end{itemize}
      \end{block}
    \end{column}
    \begin{column}{0.5\textwidth}
      \onslide<2->{\begin{block}{Sensurveiledning}
        \begin{itemize}
          \item {Sensurveiledning for mappen publiseres i løpet av semesteret og referansegruppen får anledning til å kommentere og påvirke innholdet}
        \end{itemize}
      \end{block}}
        \onslide<3->{\textbf{En del arbeid gjenstår i planlegging av mappevurderng - nærmere informasjon kommer snart}}
    \end{column}
  \end{columns}
\end{frame}

\begin{frame}{Dagens plan:}
  \begin{itemize}
    \item {Installere Anaconda}
    \item {Se på variabler og datatyper i jupyter-notebook}
  \end{itemize}
\end{frame}

%ntnutitlepage
   
%\begin{frame}{Innhold}
%   \begin{multicols}{2}
%   \tableofcontents[hideallsubsections]
%   \end{multicols}
%\end{frame}

%\AtBeginSection[]
%{
%\begin{frame}{Innhold}
%   \begin{multicols}{2}
%   \tableofcontents[currentsection, hideothersubsections]
%   \end{multicols}
%\end{frame}
%}
%

%\begin{frame}{Tema og læringsutbytte}
%   \begin{columns}
%      \begin{column}{0.5\textwidth}
%         \begin{block}{Tema: Python start}
%            \begin{enumerate}[$\bullet$]
%                  \item Installasjon:
%                     \begin{itemize}
%                        \item {Anaconda}
%                        \item {Jupyter notebook}
%                        \item {Pycharm}
%                        \item {pip for ekstra bibliotek}
%                     \end{itemize}
%                  \item Variabler
%                  \item Regneartene
%                  \item Datatyper: Heltall (\mintinline{python}{int}), flyttall (\mintinline{python}{float}) og strenger(\mintinline{python}{str})
%            \end{enumerate}
%         \end{block}
%      \end{column}
%      \begin{column}{0.5\textwidth}
%         Du skal kunne:
%         \small
%         \begin{enumerate}[$\bullet$]
%               \item Strenger
%               \item Variabler
%               \item Regneartene
%               \item print .m formatering
%               \item {importere (ikke bruke) bibliotek?}
%         \end{enumerate}
%      \begin{block}{Relevante eksempler}
%         ?? Renter osv?
%      \end{block}
%      \end{column}
%   \end{columns}
%\end{frame}

\end{document}

